\documentclass[11pt]{article}
\usepackage[margin=1in]{geometry}
\usepackage{amsmath,amssymb,amsthm}
\usepackage[T1]{fontenc}
\usepackage{lmodern}
\usepackage{hyperref}
\hypersetup{colorlinks=true,linkcolor=blue,citecolor=blue,urlcolor=blue,
  pdftitle={AIM 340 and the Alman--Vassilevska Williams 3SUM algorithm},
  pdfauthor={Kaiwen Jin}}
\newtheorem{lemma}{Lemma}
\newtheorem{proposition}[lemma]{Proposition}
\newtheorem{theorem}[lemma]{Theorem}
\newcommand{\bits}{\operatorname{bit}}
\title{AIM 340 and the Alman--Vassilevska Williams 3SUM algorithm}
\author{Kaiwen Jin\\\small\href{mailto:yui.ui.siki@gmail.com}{yui.ui.siki@gmail.com}}
\date{October 12, 2026}

\begin{document}
\maketitle

\begin{abstract}
Alman and Vassilevska Williams give a deterministic
$O(N^{1.9992})$-time algorithm for integer 3SUM in the logarithmic-word
RAM. We explain its application to the exact version in AIM Problem 340:
given $n$ distinct integers in $[-n^3,n^3]$, decide whether three pairwise
distinct elements satisfy $a+b=c$. An elementary reduction uses
$O(\log n)$ zero-sum instances of size at most $2n$, with a fixed
polynomial bound on their entries. It gives total time
$O(n^{2-1/2500})$ for every valid input.
The fast algorithm is entirely due to Alman and Vassilevska Williams;
this note records the correspondence with the stated target.
\end{abstract}

\section{The target and the cited algorithm}\label{sec:target}

We use Problem 340 of the AIM collection at the exact revision
recorded in reference~\cite{aim}.
For each integer $n\geq3$, its input is an explicitly listed set
$S\subseteq[-n^3,n^3]\cap\mathbb Z$ with $|S|=n$.
The answer is YES precisely when there are pairwise distinct
$a,b,c\in S$ with $a+b=c$.
The model is a classical word-RAM with $w=\Theta(\log n)$-bit words,
large enough to store an input integer and an address.
Usual word arithmetic, bit operations, comparisons and indexed memory
accesses have unit cost. A single finite program must work at every size.
All input-dependent preprocessing is counted.

For an admissible algorithm $A$, the target uses
\[
 T_A(n)=
 \max_{\substack{S\subseteq[-n^3,n^3]\cap\mathbb Z\\|S|=n}}
 \mathbb E[\text{total word-RAM steps of }A(S)].
\]
The answer must be correct with probability at least $2/3$ on each input.
The conjecture excludes $T_A(n)=O(n^{2-\varepsilon})$ for every fixed
$\varepsilon>0$.

The primary result we use is Alman and Vassilevska
Williams~\cite[Theorem 22, Section 3.4]{avw}.
For every fixed constant $\nu$, they give a deterministic algorithm for
zero-sum 3SUM on $N$ integers of absolute value at most $N^\nu$,
with total time $O(N^{1.9992})$.
Their Theorem 2 specifies the standard word-RAM with $O(\log N)$-bit
words. We fix $\nu=6$ throughout.
We denote this algorithm by $Z$.
Its construction, preprocessing and running-time analysis belong to
the cited work. The task here is to enforce distinctness of the
original elements while using $Z$.

\section{A reduction preserving distinct elements}\label{sec:reduction}

First sort $S$. If $0\in S$, test whether both $x$ and $-x$ belong to
$S$ for some nonzero $x$. If so, return YES:
$(x,-x,0)$ is a valid distinct triple.
Sorting and binary searches take $O(n\log n)$ total time.

Set
\[
 S^*=S\setminus\{0\},\qquad m=|S^*|,\qquad U=n^3,\qquad Q=4U+1.
\]
Assign the elements of $S^*$ their sorted ranks
$j=0,\ldots,m-1$. For $t\geq0$, write
$\bits_t(j)=\lfloor j/2^t\rfloor\bmod2$.
Let $L=\lceil\log_2m\rceil$; since $n\geq3$, we have $m\geq2$.
For $0\leq t<L$, form
\begin{align}
 T_t={}&
 \{Q+x:x\in S^*,\ \bits_t(\operatorname{rank}(x))=0\}\notag\\
 &{}\cup
 \{2Q+x:x\in S^*,\ \bits_t(\operatorname{rank}(x))=1\}\notag\\
 &{}\cup\{-3Q-x:x\in S^*\}.\label{eq:encoding}
\end{align}
Call $Z$ on each $T_t$ and return YES if any call does.
If the initial stage and all calls return NO, return NO.
The tags $1,2,-3$ force a zero-sum triple to use one value from each
of the three displayed sets. The rank bit then separates the two summands.

\begin{lemma}\label{lem:legal}
Every $T_t$ has $N=2m$ distinct integers, where $n\leq N\leq2n$,
and all its entries have absolute value at most $N^6$.
\end{lemma}
\begin{proof}
The three sets in~\eqref{eq:encoding} lie respectively in
\[
 [3U+1,5U+1],\qquad [7U+2,9U+2],\qquad
 [-13U-3,-11U-3].
\]
These intervals are disjoint, and each encoding within one set is
injective. The first two sets together have $m$ entries, and the
last has $m$. Also $n-1\leq m\leq n$, whence $n\leq2m\leq2n$.
Finally, for $n\geq3$,
\[
 13n^3+3\leq14n^3\leq n^6\leq N^6.
\]
Thus each call has the fixed polynomial integer bound required by $Z$.
\end{proof}

\begin{proposition}\label{prop:correct}
The reduction returns YES if and only if $S$ contains pairwise distinct
$a,b,c$ with $a+b=c$.
\end{proposition}
\begin{proof}
The initial opposite-pair test produces a valid witness.
Suppose a call on $T_t$ finds a zero-sum triple.
Each term has the form $h_iQ+r_i$, with
$h_i\in\{1,2,-3\}$ and $|r_i|\leq U$.
Since $|r_1+r_2+r_3|\leq3U<Q$, zero sum forces
$h_1+h_2+h_3=0$ and $r_1+r_2+r_3=0$.

Let $k_1,k_2,k_3$ count the tags $1,2,-3$.
Then $k_1+k_2+k_3=3$ and $k_1+2k_2-3k_3=0$.
If $k_3=0$ the latter sum is positive; if $k_3\geq2$ it is negative.
Hence $k_3=1$ and $k_1=k_2=1$.
The triple is therefore
$(Q+a,2Q+b,-3Q-c)$ for $a,b,c\in S^*$ with $a+b=c$.
The rank-bit condition gives $a\ne b$.
If $c=a$, then $b=0$, and if $c=b$, then $a=0$.
Both are excluded by the definition of $S^*$.
Thus all three original elements are distinct.
The same tag argument shows that the encoded values are distinct,
so allowing repeated values in a zero-sum convention would not alter
any oracle answer here.

Conversely, suppose a valid triple satisfies $a+b=c$.
When $c=0$, the initial test finds its opposite pair.
Otherwise neither summand is zero: a zero summand would make the other
equal to $c$. Thus $a,b,c\in S^*$.
The distinct ranks of $a$ and $b$ differ in some bit $t<L$.
Interchange $a,b$ if needed so that their bits are $0,1$.
Then the three corresponding values in~\eqref{eq:encoding}
sum to zero. The call on this $T_t$ returns YES.
\end{proof}

\section{Total time and the exact consequence}\label{sec:time}

Each set $T_t$ can be constructed in $O(n)$ word operations.
Its integers have $O(\log n)$ bits, and
$\log N=\Theta(\log n)$ because $n\leq N\leq2n$.
Thus the cited algorithm and the new arithmetic use the same
logarithmic-word model, with a fixed sufficiently large constant
in the word size.
The source algorithm is fixed for $\nu=6$; sorting, rank assignment,
bit tests and the loop over $t$ are also one finite program.
No advice or supplied data structure is used.

Including all preprocessing and all $O(\log n)$ source calls, the time is
\[
 O(n\log n+n^{1.9992}\log n).
\]
All logarithms in this time bound may be taken to base two.
Put $q=2500$, $k=\lfloor\log_2n\rfloor$ and
$j=\lfloor k/q\rfloor$. Then $2^{qj}\leq n$ and
$k+1\leq q(j+1)$. The elementary induction
$2^j\geq j+1$ for integers $j\geq0$ gives
\[
 \log_2n<k+1\leq q(j+1)\leq q2^j\leq qn^{1/q}.
\]
Consequently the total time is
$O(n^{1.9996})=O(n^{2-1/2500})$, uniformly over $S$.

\begin{theorem}\label{thm:main}
There is one uniform deterministic classical word-RAM algorithm $A$,
with word size $\Theta(\log n)$, which, for every integer $n\geq3$
and every set $S\subseteq[-n^3,n^3]\cap\mathbb Z$ of cardinality $n$,
correctly decides whether pairwise distinct $a,b,c\in S$ satisfy $a+b=c$.
Counting all input-dependent preprocessing, it satisfies
\[
 T_A(n)=O(n^{2-1/2500}).
\]
\end{theorem}

Here $T_A$ is the maximum expected step count defined in
Section~\ref{sec:target}; determinism makes the expectation equal to the
actual step count. The success probability on every input is $1$,
and the positive exponent improvement $\varepsilon=1/2500$ is fixed.
The constants do not depend on $S$ or $n$.
The theorem follows from the cited algorithm through
Lemma~\ref{lem:legal} and Proposition~\ref{prop:correct}, and
refutes the precise conjecture in the pinned AIM statement.
The source of the fixed-power speedup is~\cite{avw}.

\begin{thebibliography}{2}
\raggedright

\bibitem{avw}
Josh Alman and Virginia Vassilevska Williams,
\emph{Truly Subquadratic 3SUM and Truly Subcubic APSP via Triangles
in Sparse Lopsided Graphs},
arXiv:2610.06783v1, October 5, 2026.
\url{https://arxiv.org/abs/2610.06783v1}.
Theorems 2 and 22; Section 3.4.

\bibitem{aim}
Steve Brunton, Matthew J. Colbrook, Maarten V. de Hoop,
George Stepaniants, Alex Townsend and Rachel Ward,
\emph{AIM---Open Applied Problems}, Problem 340,
\emph{The modern integer 3SUM conjecture},
revision c9929805cece044705abaaebe1cebddac1ba11a6.
\url{https://github.com/MColbrook/AIM/blob/c9929805cece044705abaaebe1cebddac1ba11a6/problems/340-integer-three-sum-hardness.md}.

\end{thebibliography}

\paragraph{AI use.}
OpenAI Codex assisted the source search, reduction, checking and drafting
of this note under the author's direction.
The fast 3SUM algorithm is due entirely to Alman and Vassilevska Williams.
Independent mathematical review is requested.

\end{document}
